% =====================================================================================================================
% preambule.tex — mise en page du dossier complet CT-CORA3D (XeLaTeX)
% Texte : Libertinus Serif. Titres, légendes, tableaux, en-têtes : Avenir Next (police livrée avec macOS ; si elle
% manque, Libertinus Sans la remplace). Code : Inconsolata. Mathématiques : Libertinus Math. Icônes : Font Awesome 5.
% Couleurs : bleu marine (titres, bandeaux), terre cuite (accents), une couleur par type d'encadré.
% =====================================================================================================================
\usepackage{mathtools}                                      % avant unicode-math
\usepackage{fontspec}
\usepackage[math-style=ISO,bold-style=ISO]{unicode-math}
\setmainfont{LibertinusSerif}[Extension=.otf, UprightFont=*-Regular, ItalicFont=*-Italic,
  BoldFont=*-Semibold, BoldItalicFont=*-SemiboldItalic, Scale=1.04]
\setmathfont{LibertinusMath-Regular.otf}[Scale=1.04]
\setmonofont{Inconsolatazi4}[Extension=.otf, UprightFont=*-Regular, BoldFont=*-Bold, Scale=MatchLowercase]
\IfFontExistsTF{Avenir Next}{%
  \setsansfont{Avenir Next}[UprightFont={* Regular}, ItalicFont={* Italic}, BoldFont={* Demi Bold},
    BoldItalicFont={* Demi Bold Italic}, Scale=MatchLowercase]
  \newfontfamily\policegrasse{Avenir Next}[UprightFont={* Bold}, Scale=MatchLowercase]
  \newfontfamily\policefine{Avenir Next}[UprightFont={* Ultra Light}, Scale=MatchLowercase]
  \newfontfamily\policemoyenne{Avenir Next}[UprightFont={* Medium}, BoldFont={* Demi Bold}, Scale=MatchLowercase]
}{%
  \setsansfont{LibertinusSans}[Extension=.otf, UprightFont=*-Regular, ItalicFont=*-Italic, BoldFont=*-Bold]
  \newfontfamily\policegrasse{LibertinusSans}[Extension=.otf, UprightFont=*-Bold]
  \newfontfamily\policefine{LibertinusSans}[Extension=.otf, UprightFont=*-Regular]
  \newfontfamily\policemoyenne{LibertinusSans}[Extension=.otf, UprightFont=*-Regular, BoldFont=*-Bold]
}
\usepackage[french]{babel}
\frenchsetup{ItemLabels=\textendash, StandardItemizeEnv=true}

\usepackage[a4paper, left=2.5cm, right=2.5cm, top=2.8cm, bottom=2.8cm, headheight=16pt, headsep=0.7cm,
  footskip=1.2cm]{geometry}
\usepackage{microtype}
\setlength{\parskip}{0.5em}
\setlength{\parindent}{0pt}
\linespread{1.08}
\emergencystretch=1.5em

% --------------------------------------------------------------------------------------------------------- couleurs
\usepackage[table]{xcolor}
\definecolor{titre}{HTML}{1B4F72}      % bleu marine : titres, bandeaux, liens
\definecolor{nuit}{HTML}{0F2B40}       % marine profond : couverture, pages de partie
\definecolor{halo}{HTML}{245A80}       % halo clair derrière le rendu de la couverture
\definecolor{accent}{HTML}{C2602A}     % terre cuite : numéros, filets, étiquettes
\definecolor{accentclair}{HTML}{F0A877}
\definecolor{encre}{HTML}{1E2832}      % couleur du texte
\definecolor{ardoise}{HTML}{55636F}    % texte secondaire (légendes, en-têtes)
\definecolor{filet}{HTML}{D3DBE2}      % filets fins
\definecolor{zebre}{HTML}{F3F6F9}      % lignes alternées des tableaux
\definecolor{entete}{HTML}{DFE9F2}     % ligne de titre des tableaux
\definecolor{segm}{HTML}{1F5F8B}       \definecolor{segmfond}{HTML}{EAF2FA}
\definecolor{planif}{HTML}{B8611A}     \definecolor{planiffond}{HTML}{FDF2E6}
\definecolor{guide}{HTML}{3A7D2C}      \definecolor{guidefond}{HTML}{EBF6E7}
\definecolor{bloc}{HTML}{6F4A99}       \definecolor{blocfond}{HTML}{F3EDF8}
\definecolor{alerte}{HTML}{B03A3A}     \definecolor{alertefond}{HTML}{FBEBEB}
\definecolor{gris}{HTML}{5B6B7A}       \definecolor{grisfond}{HTML}{F2F4F7}
\definecolor{code}{HTML}{2D3E50}       \definecolor{codefond}{HTML}{F5F7FA}
\makeatletter                                               % texte en « encre » plutôt qu'en noir pur
\AtBeginDocument{\color{encre}\let\default@color\current@color\let\XC@default@color\XC@current@color}
\makeatother

% --------------------------------------------------------------------------------------------------------- outils
\usepackage{graphicx}
\usepackage{float}
\usepackage{placeins}
\renewcommand{\topfraction}{0.85}          % une figure peut occuper presque toute la page
\renewcommand{\bottomfraction}{0.6}
\renewcommand{\textfraction}{0.1}
\renewcommand{\floatpagefraction}{0.75}
\setcounter{topnumber}{3}
\setcounter{totalnumber}{5}
\usepackage{caption}
\DeclareCaptionLabelSeparator{barre}{\enspace{\color{accent!45}\rule[-0.2em]{0.7pt}{1.05em}}\enspace}
\captionsetup{font={small,sf}, labelfont={sf,bf,color=accent}, textfont={color=ardoise}, labelsep=barre,
  margin=0.4cm, skip=8pt}
\captionsetup[table]{position=top, skip=6pt}
\usepackage{subcaption}
\usepackage{booktabs, tabularx, longtable, multirow, makecell, array}
\newcolumntype{L}[1]{>{\raggedright\arraybackslash}p{#1}}
\newcolumntype{Y}{>{\raggedright\arraybackslash}X}
\usepackage{enumitem}
\setlist{itemsep=0.2em, topsep=0.35em, parsep=0.1em}
\setlist[itemize,1]{label={\color{accent}\rule[0.1em]{0.33em}{0.33em}}, leftmargin=1.35em}
\setlist[itemize,2]{label={\color{titre}\textendash}, leftmargin=1.3em}
\setlist[enumerate,1]{label={\sffamily\bfseries\color{titre}\arabic*.}, leftmargin=1.7em}
\usepackage{siunitx}
\sisetup{locale=FR, group-minimum-digits=4}
\usepackage{pdflscape}
\usepackage{tikz}
\usetikzlibrary{calc, positioning, fadings}
\usepackage[most]{tcolorbox}
\usepackage{listings}
\usepackage{xurl}
\usepackage{fontawesome5}

% --------------------------------------------------------------------------------------------------------- tableaux
% Ligne de titre teintée (ajoutée par \toprule), lignes alternées, filets bleu marine, pas d'espace blanc autour
% des filets (sinon les lignes colorées seraient coupées).
\renewcommand{\arraystretch}{1.32}
\setlength{\aboverulesep}{0pt}
\setlength{\belowrulesep}{0pt}
\setlength{\heavyrulewidth}{0.9pt}
\setlength{\lightrulewidth}{0.5pt}
\arrayrulecolor{titre}
\let\topruleorigine\toprule
\renewcommand{\toprule}{\topruleorigine\rowcolor{entete}}
\newcommand{\titrecol}[1]{{\sffamily\bfseries\color{titre}#1}}   % texte de la ligne de titre
\AtBeginDocument{\rowcolors{2}{zebre}{white}}

% --------------------------------------------------------------------------------------------------------- titres
\usepackage[explicit]{titlesec}
\newcommand{\espacelettres}[1]{{\addfontfeature{LetterSpace=14}#1}}          % lettres espacées
\newcommand{\etiquette}[1]{{\sffamily\bfseries\addfontfeature{LetterSpace=16}\MakeUppercase{#1}}}
\newcommand{\etiquettesansnumero}{CT-CORA3D · dossier complet}
\newif\ifnumerochapitre
\newif\ifannexe
\newlength{\hauteurbandeau}
\setlength{\hauteurbandeau}{6.5cm}
% bandeau d'ouverture de chapitre : #1 étiquette (« Chapitre 6 »), #2 numéro décoratif, #3 titre
\newcommand{\bandeau}[3]{%
  \begin{tikzpicture}[remember picture, overlay]
    \coordinate (bhg) at (current page.north west);
    \coordinate (bbd) at ([yshift=-\hauteurbandeau]current page.north east);
    \fill[titre] (bhg) rectangle (bbd);
    \begin{scope}
      \clip (bhg) rectangle (bbd);
      \node[anchor=south east, inner sep=0pt, text=white, opacity=0.10,
        font=\policegrasse\fontsize{200}{200}\selectfont] at ([xshift=-1.0cm, yshift=-0.9cm]bbd) {#2};
    \end{scope}
    \fill[accent] ([yshift=-\hauteurbandeau]current page.north west) rectangle ([yshift=-2.6pt]bbd);
    \node[anchor=south west, inner sep=0pt, text width=13.6cm, text=white]   % (align=left forcerait des césures)
      at ([xshift=2.5cm, yshift=-\hauteurbandeau+1.25cm]current page.north west)
      {\raggedright{\fontsize{8.6}{11}\selectfont\color{accentclair}\etiquette{#1}}\par\vspace{9pt}%
       {\sffamily\bfseries\fontsize{25}{30.5}\selectfont #3\par}};
  \end{tikzpicture}}
\titleformat{\chapter}[block]{}{}{0pt}
  {\global\numerochapitretrue\bandeau{\chaptertitlename~\thechapter}{\thechapter}{#1}}
\titleformat{name=\chapter, numberless}[block]{}{}{0pt}
  {\global\numerochapitrefalse\bandeau{\etiquettesansnumero}{}{#1}}
\newlength{\espacebandeau}                 % espace entre le haut du texte et la fin du bandeau
\setlength{\espacebandeau}{\dimexpr\hauteurbandeau-1.95cm\relax}
\titlespacing*{\chapter}{0pt}{-\topskip}{\espacebandeau}

\newcommand{\formatsection}{\sffamily\bfseries\fontsize{15}{18.5}\selectfont\color{titre}\raggedright}
\newcommand{\formatsoussection}{\sffamily\bfseries\fontsize{12.3}{15.2}\selectfont\color{titre}\raggedright}
\titleformat{\section}{\formatsection}{\textcolor{accent}{\thesection}}{0.75em}{#1}
\titleformat{name=\section, numberless}{\formatsection}{}{0pt}{#1}
\titleformat{\subsection}{\formatsoussection}{\textcolor{accent}{\thesubsection}}{0.7em}{#1}
\titleformat{name=\subsection, numberless}{\formatsoussection}{}{0pt}{#1}
\titleformat{\subsubsection}{\sffamily\bfseries\normalsize\color{ardoise}\raggedright}{}{0em}{#1}
\titlespacing*{\section}{0pt}{1.7em}{0.55em}
\titlespacing*{\subsection}{0pt}{1.3em}{0.4em}
\titlespacing*{\subsubsection}{0pt}{1em}{0.25em}
\setcounter{secnumdepth}{2}
\setcounter{tocdepth}{1}

% premier paragraphe d'un chapitre (ce que contient le chapitre)
\newenvironment{chapeau}
  {\par\begingroup\fontsize{12.3}{17}\selectfont\color{ardoise}}
  {\par\endgroup\vspace{2pt}{\color{accent}\rule{2.4cm}{1.3pt}}\par\vspace{5pt}}

% --------------------------------------------------------------------------------------------------------- en-têtes
\usepackage{fancyhdr}
\pagestyle{fancy}
\fancyhf{}
\renewcommand{\chaptermark}[1]{\markboth{#1}{}}
\newcommand{\styleentete}{\sffamily\fontsize{8.4}{10}\selectfont}
\newcommand{\numeropage}{{\sffamily\bfseries\fontsize{9.5}{11}\selectfont\color{titre}\thepage}}
% onglet sur le bord droit de la page : sa hauteur dépend du chapitre (chapitres 1 à 14, puis annexes A à E)
\newcommand{\onglet}{%
  \ifnumerochapitre
  \begin{tikzpicture}[remember picture, overlay]
    \pgfmathsetlengthmacro{\hautonglet}{2.3cm + (\the\numexpr\ifannexe 14+\fi\value{chapter}\relax - 1)*1.33cm}
    \fill[\ifannexe accent\else titre\fi, rounded corners=3pt]
      ([xshift=-0.78cm, yshift=-\hautonglet]current page.north east) rectangle
      ([xshift=0.2cm, yshift=-\hautonglet-1.18cm]current page.north east);
    \node[text=white, font=\policegrasse\fontsize{10.5}{12}\selectfont]
      at ([xshift=-0.39cm, yshift=-\hautonglet-0.59cm]current page.north east) {\thechapter};
  \end{tikzpicture}%
  \fi}
\fancyhead[L]{\styleentete\color{ardoise}\ifnumerochapitre{\bfseries\color{accent}\thechapter}\hspace{0.7em}\fi
  \nouppercase{\leftmark}}
\fancyhead[R]{\styleentete\bfseries\color{titre}\espacelettres{CT-CORA3D}}
\fancyfoot[C]{\onglet}
\fancyfoot[R]{\numeropage}
\renewcommand{\headrulewidth}{0.4pt}
\renewcommand{\headrule}{{\color{filet}\hrule width\headwidth height\headrulewidth}}
\fancypagestyle{plain}{\fancyhf{}\fancyfoot[R]{\numeropage}\renewcommand{\headrulewidth}{0pt}}
\fancypagestyle{paysage}{\fancyhf{}\fancyfoot[R]{\numeropage}\renewcommand{\headrulewidth}{0pt}}

% --------------------------------------------------------------------------------------------------------- liens
\usepackage{hyperref}
\hypersetup{colorlinks=true, linkcolor=titre, citecolor=accent, urlcolor=segm,
  pdftitle={CT-CORA3D — Dossier complet du projet}, pdfauthor={Dorian Verstappen},
  pdfsubject={Planification CORA 3D et guides de coupe imprimables chez le chien}, pdflang={fr-FR}}
\usepackage{bookmark}
\usepackage[french, noabbrev, nameinlink]{cleveref}
\crefname{figure}{figure}{figures}            \Crefname{figure}{Figure}{Figures}
\crefname{table}{tableau}{tableaux}           \Crefname{table}{Tableau}{Tableaux}
\crefname{equation}{équation}{équations}      \Crefname{equation}{Équation}{Équations}
\crefname{chapter}{chapitre}{chapitres}       \Crefname{chapter}{Chapitre}{Chapitres}
\crefname{section}{section}{sections}         \Crefname{section}{Section}{Sections}
\crefname{subsection}{section}{sections}      \Crefname{subsection}{Section}{Sections}
\crefname{appendix}{annexe}{annexes}          \Crefname{appendix}{Annexe}{Annexes}
\creflabelformat{equation}{#2(#1)#3}

\numberwithin{equation}{chapter}

% --------------------------------------------------------------------------------------------------------- sommaire
\usepackage{etoc}
\newcommand{\pointstoc}{\leaders\hbox to 0.75em{\hss\textcolor{ardoise!55}{.}\hss}\hfill}
\newcommand{\lignetoc}[6][\hfill]{% [#1 remplissage], #2 retrait, #3 largeur du numéro, #4 numéro, #5 titre, #6 page
  {\leftskip=#2\relax \rightskip=1.3cm plus 1fil\relax \parfillskip=-\rightskip \parindent=0pt
   \noindent\llap{\makebox[#3][l]{#4}}#5\nobreak#1\nobreak\makebox[1.3cm][r]{#6}\par}}
\newcommand{\tabledesmatieres}{%
  \begingroup
  \etocsettocstyle{\chapter*{\contentsname}\markboth{\contentsname}{}}{}%
  \etocsetstyle{part}{}{\addvspace{15pt}}
    {\begingroup\hypersetup{linkcolor=white}\noindent
     \colorbox{titre}{\makebox[\dimexpr\linewidth-2\fboxsep\relax][l]{\rule[-0.6em]{0pt}{1.85em}\hspace{0.45em}%
       \etocifnumbered{{\sffamily\bfseries\fontsize{8.4}{10}\selectfont\color{accentclair}%
         \espacelettres{PARTIE~\etocthenumber}}\hspace{1.1em}}{}%
       {\sffamily\bfseries\color{white}\etocname}\hfill
       {\sffamily\bfseries\color{white}\etocpage}\hspace{0.45em}}}\par\endgroup\nobreak\addvspace{3pt}}{}%
  \etocsetstyle{chapter}{}{\addvspace{8pt}}
    {\begingroup\hypersetup{linkcolor=titre}%
     \lignetoc{1.25cm}{1.25cm}{\etocifnumbered{\policegrasse\fontsize{12.5}{14}\selectfont\color{accent}%
       \etocthenumber}{}}{\sffamily\bfseries\fontsize{11.3}{14}\selectfont\etocname}%
       {\sffamily\bfseries\color{titre}\etocpage}\endgroup\nobreak\addvspace{1pt}}{}%
  \etocsetstyle{section}{}{}
    {\begingroup\hypersetup{linkcolor=encre}%
     \lignetoc[\pointstoc]{2.35cm}{1.1cm}{\sffamily\fontsize{8.8}{11}\selectfont\color{ardoise}\etocthenumber}%
       {\fontsize{10.4}{12.8}\selectfont\etocname}%
       {\sffamily\fontsize{8.8}{11}\selectfont\color{ardoise}\etocpage}\endgroup}{}%
  \etocsettocdepth{section}%
  \tableofcontents
  \endgroup}

% --------------------------------------------------------------------------------------------------------- parties
% \partie{titre}{description} : page pleine en bleu nuit, grand numéro, description et liste des chapitres ;
% \partieannexes{description} : la même page pour les annexes (sans numéro de partie)
\newcommand{\pagepartie}[4]{% #1 étiquette, #2 grand signe décoratif, #3 titre, #4 description
  \begin{tikzpicture}[remember picture, overlay]
    \fill[nuit] (current page.south west) rectangle (current page.north east);
    \begin{scope}
      \clip (current page.south west) rectangle (current page.north east);
      \node[anchor=south east, inner sep=0pt, text=white, opacity=0.07,
        font=\policegrasse\fontsize{360}{360}\selectfont] at ([xshift=-0.8cm, yshift=-1.6cm]current page.south east)
        {#2};
    \end{scope}
  \end{tikzpicture}\par
  \vspace*{\dimexpr3.3cm-\baselineskip\relax}%
  \begingroup
  \raggedright\color{white}\hyphenpenalty=10000
  {\fontsize{9.5}{11}\selectfont\color{accentclair}\etiquette{#1}\par}%
  \vspace{12pt}%
  {\sffamily\bfseries\fontsize{29}{35}\selectfont #3\par}%
  \vspace{16pt}%
  {\color{accent}\rule{2.2cm}{2.6pt}\par}%
  \vspace{16pt}%
  {\sffamily\fontsize{12.6}{18}\selectfont\color{white!78!nuit}#4\par}%
  \vspace{1.3cm}%
  \hypersetup{linkcolor=white}%
  \etocsetstyle{chapter}{}{\addvspace{7pt}}
    {\lignetoc{1.3cm}{1.3cm}{\policegrasse\fontsize{13}{15}\selectfont\color{accentclair}\etocthenumber}%
       {\sffamily\fontsize{12.4}{15.5}\selectfont\color{white}\etocname}%
       {\sffamily\fontsize{11}{13}\selectfont\color{white!60!nuit}\etocthepage}}{}%
  \etocsettocstyle{}{}%
  \etocsetnexttocdepth{chapter}%
  \localtableofcontents
  \endgroup
  \clearpage}
\newcommand{\partie}[2]{%
  \clearpage
  \thispagestyle{empty}%
  \refstepcounter{part}%
  \addcontentsline{toc}{part}{\protect\numberline{\thepart}#1}%
  \global\numerochapitrefalse
  \pagepartie{Partie \Roman{part}}{\Roman{part}}{#1}{#2}}
\newcommand{\partieannexes}[1]{%
  \clearpage
  \thispagestyle{empty}%
  \phantomsection
  \addcontentsline{toc}{part}{Annexes}%
  \global\numerochapitrefalse
  \pagepartie{Pour aller plus loin}{A}{Annexes}{#1}}

% --------------------------------------------------------------------------------------------------------- encadrés
% Fond teinté, filet de couleur à gauche, pastille ronde avec une icône posée sur le filet.
\newcommand{\pastille}[2]{\node[circle, fill=#1, text=white, inner sep=0pt, minimum size=17pt,
  font=\fontsize{8.3}{8.3}\selectfont] at ([xshift=1.5pt, yshift=-11pt]frame.north west) {#2};}
\tcbset{
  boite/.style n args={3}{enhanced, sharp corners=west, arc=3pt, boxrule=0pt, leftrule=3pt,
    colframe=#1, colback=#2, colupper=encre, left=14pt, right=10pt, top=6pt, bottom=7pt,
    before skip=0.95em, after skip=0.95em, fonttitle=\sffamily\bfseries\small, coltitle=#1,
    attach title to upper={\par\smallskip}, parbox=false, overlay={\pastille{#1}{#3}}},
  highlight math style={enhanced, colback=accent!5, colframe=accent!55, boxrule=0.6pt, arc=3pt,
    left=7pt, right=7pt, top=4pt, bottom=4pt}}
\newtcolorbox{aretenir}[1][]{boite={segm}{segmfond}{\faLightbulb[regular]}, title={À retenir}, #1}
\newtcolorbox{danslecode}[1][]{boite={gris}{grisfond}{\faCode}, title={Dans le code}, fontupper=\small, #1}
\newtcolorbox{exemple}[1][]{boite={guide}{guidefond}{\faCalculator}, title={Exemple chiffré}, #1}
\newtcolorbox{attention}[1][]{boite={planif}{planiffond}{\faExclamationTriangle}, title={Attention}, #1}
\newtcolorbox{pourquoi}[1][]{boite={bloc}{blocfond}{\faQuestion}, title={Pourquoi ce choix ?}, #1}
% carte du guide de lecture : #1 icône, #2 lecteur
\newtcolorbox{carte}[2]{enhanced, colback=white, colframe=filet, boxrule=0.7pt, arc=4pt, left=11pt, right=11pt,
  top=10pt, bottom=10pt, fontupper=\small, colupper=encre, title={#2}, fonttitle=\sffamily\bfseries\normalsize,
  coltitle=titre, attach title to upper={\par\smallskip},
  overlay={\node[anchor=north east, text=accent!80, font=\fontsize{17}{17}\selectfont]
    at ([xshift=-10pt, yshift=-9pt]frame.north east) {#1};}}

% --------------------------------------------------------------------------------------------------------- code
\lstdefinestyle{dossier}{basicstyle=\ttfamily\small\color{code}, keywordstyle=\color{segm}\bfseries,
  commentstyle=\color{gris!90}, stringstyle=\color{planif}, numbers=none, frame=l, framerule=2.5pt,
  rulecolor=\color{titre!40}, backgroundcolor=\color{codefond}, xleftmargin=10pt, framexleftmargin=8pt,
  framextopmargin=3pt, framexbottommargin=3pt, aboveskip=0.8em, belowskip=0.8em, showstringspaces=false,
  columns=fullflexible, keepspaces=true, breaklines=true, upquote=true,
  literate={é}{{\'e}}1 {è}{{\`e}}1 {ê}{{\^e}}1 {à}{{\`a}}1 {â}{{\^a}}1 {ù}{{\`u}}1 {û}{{\^u}}1 {ô}{{\^o}}1
           {î}{{\^i}}1 {ç}{{\c c}}1 {É}{{\'E}}1 {À}{{\`A}}1 {θ}{{$\theta$}}1 {→}{{$\to$}}1 {−}{{-}}1}
\lstset{style=dossier}

% --------------------------------------------------------------------------------------------------------- macros
\newcommand{\fichier}[1]{\nolinkurl{#1}}                  % nom de fichier ou de dossier
\ExplSyntaxOn                                              % nom de fonction, de clé, de variable :
\tl_new:N \l__dossier_code_tl                              % coupure de ligne permise après chaque « _ »
\NewDocumentCommand{\code}{m}
  {
    \tl_set:Nx \l__dossier_code_tl { \tl_to_str:n {#1} }
    \regex_replace_all:nnN { \_ } { \c{textunderscore} \c{allowbreak} } \l__dossier_code_tl
    \texttt { \tl_use:N \l__dossier_code_tl }
  }
\ExplSyntaxOff
\newcommand{\vb}[1]{\symbf{#1}}                            % vecteur
\newcommand{\mt}[1]{\symbf{#1}}                            % matrice
\DeclarePairedDelimiter{\norm}{\lVert}{\rVert}
\DeclarePairedDelimiter{\abs}{\lvert}{\rvert}
\DeclareMathOperator{\atantwo}{atan2}
\DeclareMathOperator{\clamp}{clamp}
\DeclareMathOperator{\smin}{smin}
\newcommand{\dd}{\,\mathrm{d}}
\newcommand{\RR}{\mathbb{R}}
\renewcommand{\degre}{\ensuremath{^\circ}}
\newcommand{\etape}[1]{\textsf{\bfseries\color{titre}#1}}  % nom d'étape dans le texte
\newcommand{\ui}[1]{\textsf{«\,#1\,»}}                      % libellé de l'interface
\newcommand{\remarque}[1]{{\small\color{gris}#1}}

% figure sur toute la largeur, chemin relatif au dossier figures/
\graphicspath{{figures/}{figures/externes/}}
\newcommand{\fig}[4][1]{%
  \begin{figure}[htbp]\centering\includegraphics[width=#1\linewidth]{#2}\caption{#3}\label{#4}\end{figure}}
% une image d'une planche de rendus 3D : #1 hauteur, #2 image, #3 lettre, #4 légende courte
\newsavebox{\boitepanneau}
\newcommand{\panneau}[4]{%
  \sbox{\boitepanneau}{\includegraphics[height=#1]{#2}}%
  \begin{minipage}[t]{\ifdim\wd\boitepanneau<2.9cm 2.9cm\else\wd\boitepanneau\fi}\centering
    \usebox{\boitepanneau}\par\vspace{5pt}
    {\sffamily\footnotesize\color{ardoise}{\bfseries\color{accent}#3}\enspace #4\par}%
  \end{minipage}}
